\chapter{Criterios de desempeño en sistemas de control}
\label{ch:criterios-desempeno}

\section{Introducción del capítulo}
No basta con afirmar que un controlador ``funciona''; es necesario establecer cómo se mide su
calidad. En problemas de regulación de frecuencia, la percepción de buen desempeño depende
de variables concretas: qué tan profunda es la desviación inicial, cuánto tarda el sistema en
recuperarse y cuánta energía acumulada queda asociada al error. Estas métricas permiten pasar
de una evaluación cualitativa a una comparación técnicamente defendible.

\section{Objetivo del capítulo}
El propósito de este capítulo es presentar los indicadores más utilizados para valorar la
respuesta dinámica de un sistema controlado y justificar cuáles de ellos resultan más útiles
para una comparativa entre sintonización clásica y sintonización optimizada. Se busca, además,
conectar cada criterio con implicancias físicas reales sobre calidad de energía, esfuerzo del
actuador y estabilidad global.

\section{Indicadores temporales}
Los indicadores temporales describen la forma visible de la respuesta en el dominio del tiempo.
El tiempo de subida informa cuán rápido reacciona el sistema; el sobreimpulso máximo muestra
la severidad del pico transitorio; el tiempo de asentamiento refleja cuánto demora la variable en
permanecer dentro de una banda aceptable; y el error en estado estacionario evidencia si el
controlador logra o no eliminar la desviación residual. En microredes aisladas, estos
indicadores se vinculan directamente con la continuidad del servicio y la protección de equipos.

\section{Índices integrales de error}
Los índices integrales condensan el desempeño completo de la respuesta en un único valor
acumulado. El IAE suma el valor absoluto del error; el ISE penaliza con mayor intensidad los
errores grandes; el ITSE enfatiza errores tardíos al multiplicarlos por el tiempo y por su valor
cuadrático; y el ITAE pondera el error absoluto con el tiempo, castigando de manera especial
las oscilaciones persistentes. Por esta razón, el ITAE es especialmente atractivo cuando se
busca una respuesta rápida, limpia y con mínima cola transitoria.

\section{Robustez y sensibilidad}
Un controlador no debe evaluarse solo bajo condiciones ideales, sino también por su capacidad
de sostener el desempeño cuando cambian la carga, la irradiancia o ciertos parámetros de la
planta. La robustez se refiere precisamente a esa capacidad de conservar estabilidad y calidad
de respuesta frente a incertidumbres y perturbaciones. La sensibilidad, por su parte, ayuda a
identificar cuánto se altera la salida ante variaciones del modelo o del lazo de control; un
sistema excesivamente sensible suele ser frágil en operación real.

\section{Criterios para la comparación de sintonías}
Para el enfoque de este libro se emplearán tres niveles de comparación. En primer lugar, se
considerarán sobreimpulso, tiempo de asentamiento y error estacionario como indicadores
directamente interpretables desde la operación del sistema. En segundo lugar, se utilizará el
ITAE como función objetivo principal porque resume la penalización sobre errores persistentes.
Finalmente, se incorporará una lectura cualitativa de robustez, asociada al esfuerzo del
gobernador y a la suavidad de la recuperación de frecuencia.

\section{Cierre del capítulo}
Una vez definidos los criterios de evaluación, es posible revisar con mayor precisión los
métodos tradicionales de ajuste y entender por qué, en ciertos sistemas, estos ya no resultan
suficientes. El capítulo siguiente examina precisamente la sintonización clásica del PID y su
alcance real frente a plantas complejas como las microredes hidro-solares aisladas.
